% =====================================================================
\section{Limitations and Future Work}\label{sec:limitations}
\emph{Validator strength.} Validators are the principal semantic boundary. Engineers must show that observations include all the information the desired relation needs, including negative dependencies and shared context, and parsers and finite test suites usually establish weaker relations. Theorem~\ref{thm:table} discharges the obligation for one finite language by making the validator the specification. RQ6 measures what happens otherwise: on standard benchmarks, a sizable fraction of base-test-accepted solutions fail stronger tests. The coordination guarantees are unaffected, but the released artifacts are only as good as the checks. Property-based validators against abstract models (RQ5), mutation-adequate test validators and proof-carrying proposals checked by a trusted verifier are routes to stronger relations.

\emph{Structural rigidity.} Agents fill slots; they cannot create objects, modules or references. This suits workflows whose structure is declared (Section~\ref{sec:intro}), but not open-ended development in which implementing a feature reveals the need for a new helper module or endpoint. Such structural change can be admitted through the structured-ingress extension (Section~\ref{sec:extensions}): a planning agent proposes a typed source delta, for example ``add endpoint \texttt{restore(id)} serving request~19''. The controller checks the delta as an untrusted Edit, and normalization then creates the new slots. This keeps every guarantee, but someone must define the delta language and its invariants, and fidelity to intent needs a checker or a human. Automatically deriving rules and read sets from code, for example from import and call graphs, and recording the inputs a prompt renderer actually reads, as dynamic build systems do~\cite{erdweg2015pluto}, are the most promising ways to reduce the declaration effort of Table~\ref{tab:roles}.

\emph{Cycles.} The acyclic discipline excludes unrestricted collaborative editing and cyclic consistency restoration. Bounded feedback epochs (Proposition~\ref{prop:epochs}) admit test-and-fix loops when a deterministic promotion policy and an epoch bound are given. They guarantee termination, not convergence to a correct result.

\emph{Early cutoff under stochastic generation.} Early cutoff requires regenerated outputs to be \emph{equal}. RQ5 and RQ6 show that this is frequent for narrowly typed slots and rare for free-form code. Normalizing slot values (canonical JSON, structure-only normal forms of code) raises the rate, and the normal form must then be part of the contract. Semantic equivalence checks would raise it further, but they reintroduce the validator-strength problem.

\emph{Trusted computing base and effects.} The guarantees assume correct normalization, key construction, authorization, transaction isolation, validator execution and publication. A production controller should be shown to refine the transitions of Table~\ref{tab:rules}, including durable request identities and crash recovery; ours keep identities for one process lifetime. Generated code is treated as inert data and executed only in a sandbox for validation. An agent with its own file-system, network or deployment permissions can cause effects outside the controller, which require mediated tool actions and a larger transition system.

\emph{Mechanization.} The main remaining verification task is a parameterized TLAPS proof of the inductive invariant, the history-counting induction and the refinement mapping, which are currently checked by TLC on bounded instances, together with connecting the structural kernel to the protocol model. Section~\ref{sec:artifact} lists the technical obstacles.

% =====================================================================
\section{Conclusion}\label{sec:conclusion}
Reliable collaboration among generative agents is a coordination problem as much as a generation problem, and coordination is where MAS research offers the right starting point: an objective, artifact-mediated medium that enforces the rules outside the agents. Generated artifacts add three requirements to this lineage: content currentness relative to what an agent was shown, exact and minimal invalidation, and bounded effort per dependency episode. Exact dependency keys and an optimistic commit guard exclude stale, unauthorized and inconsistent results; an upstream gate makes repair cost independent of scheduling; a publication barrier separates local acceptance from global readiness; and transactional reconfiguration extends these properties to changes of team and rules. The evaluation confirms that each mechanism is necessary, that exactness pays off on independent and coupled workspaces, and that the guarantees hold with real LLMs. It also shows the boundary: with real models, late replies to changed questions are often wrong and are always rejected, but correctness of what is published never exceeds the validators that admitted it.
